\documentclass[11pt]{article}
\usepackage[margin=1in]{geometry}
\usepackage[T1]{fontenc}
\usepackage[utf8]{inputenc}
\usepackage{lmodern,microtype,amsmath,amssymb,amsthm,mathtools}
\usepackage[colorlinks=true,linkcolor=blue,citecolor=blue,urlcolor=blue]{hyperref}
\hypersetup{pdftitle={Exponential online mistakes at every width for isotropic threshold networks with a residual},pdfauthor={ChatGPT-6 Astra; Kaarel Hanni; Dmitry Vaintrob}}
\newtheorem{theorem}{Theorem}
\newtheorem{lemma}[theorem]{Lemma}
\newtheorem{proposition}[theorem]{Proposition}
\newtheorem{corollary}[theorem]{Corollary}
\newtheorem{remark}[theorem]{Remark}
\newcommand{\R}{\mathbb R}
\newcommand{\E}{\mathbb E}
\newcommand{\Pp}{\mathbb P}
\newcommand{\sgn}{\operatorname{sign}}
\newcommand{\phiR}{\phi}
\title{Exponential online mistakes at every width for\protect\\isotropic threshold networks with a residual}
\author{ChatGPT-6 Astra \and Kaarel H\"anni \and Dmitry Vaintrob\thanks{Authors are listed in alphabetical order.}}
\date{6 November 2026}
\begin{document}
\maketitle

\begin{abstract}
We construct a Bayesian neural-network learner that has polynomial IID
sample complexity but makes exponentially many mistakes on a fixed
input order, at every hidden width in its strong-prior limit.
The activations are the fixed discontinuous functions
$\sgn(t)+\tanh(t)$ and $\sgn(t)+t$: a threshold neuron with a bounded
or linear residual. The bounded version also works with an arbitrarily
small fixed positive coefficient on its tanh residual.
The network has one hidden layer, the full Boolean input cube, a constant
input coordinate, and the standard independent Gaussian weight prior
with equal variances across all first-layer coordinates. The prior is
conditioned on global output margin greater than $1/2$ and exact labels.
For a target of Boolean circuit size $\Theta(k)$, a calibration prefix of
$O(k)$ examples is followed by $2^k-k-1$ designated mistakes. The target
itself has a width-one teacher with every weight bounded by one.
An exact all-competitor inequality reduces the minimum Gaussian energy,
at every width, to a strictly increasing function of the affine margin
norm. A separate convex-hull argument gives an energy gap to all
nonthreshold functions, yielding polynomial IID sample complexity.
We also prove direct width-dependent variance statements: every fixed
exponent $\gamma>1$ works at all sufficiently large widths, and a
sufficiently small mean-field prefactor works uniformly at every width.
The proof uses the activation's jump at zero; no conclusion
for ordinary tanh, ReLU, or hardtanh is claimed.
\end{abstract}

\section{The learner and the separation}

We study two fixed activations, with $\sgn(0)=0$:
\[
 \phi_{\rm lin}(t)=\sgn(t)+t,\qquad
 \phi_{\rm bnd}(t)=\sgn(t)+\tanh(t).
\]
Both are odd and have a jump at zero. The second is bounded in absolute
value by two. Each is a threshold activation with a residual, rather
than a continuous approximation to a threshold. In the model below,
$\phi$ may be either one.
For $x\in X=\{-1,1\}^d$, write $\bar x=(1,x)$.
A width-$N$ learner has output
\begin{equation}\label{eq:model}
 F_{u,v}(x)=\sum_{j=1}^N v_j\phi(u_j\cdot\bar x),
 \qquad
 u_{ji}\sim N(0,N\sigma^2/d),\quad
 v_j\sim N(0,\sigma^2/N),
\end{equation}
with all weights independent. The bias is the coefficient of the
constant input and has the same variance as every other first-layer
coordinate. There is no coordinate-dependent covariance or extra
readout regularization factor.

For a Boolean function $g$ let
\[
 U_g=\{(u,v):g(x)F_{u,v}(x)>1/2\text{ for every }x\in X\}.
\]
The posterior first conditions on the union of these full-cube
strict-margin cells and then retains exactly the cells consistent with
the observed labels. A prediction's error probability is the posterior
mass of functions assigning the wrong label. This is the error
probability of a fresh posterior sample. Posterior-label majority is
wrong when that mass exceeds $1/2$; the stronger convention requiring
correct-label probability greater than $2/3$ is also considered.

The strong-prior limit at a fixed width means $\sigma\downarrow0$.
For direct width-dependent statements we use
\begin{equation}\label{eq:scaling}
 \sigma_N^2=\tau^2N^{-\gamma},\qquad
 \operatorname{Var}(u_{ji})=\tau^2N^{1-\gamma}/d,
 \quad \operatorname{Var}(v_j)=\tau^2N^{-\gamma-1}.
\end{equation}
Here $\tau>0$ and $\gamma\ge1$ are fixed.

\begin{theorem}\label{thm:main}
For every $k\ge3$, put $d=2k+1$ and write inputs as $(x,y,z)$, where
$x,y\in\{-1,1\}^k$ and $z\in\{-1,1\}$. Define
\[
 X_x=\sum_{i=0}^{k-1}2^i\frac{x_i+1}{2},\qquad
 Y_y=\sum_{i=0}^{k-1}2^i\frac{y_i+1}{2},\qquad
 f_k(x,y,z)=+1\ \Longleftrightarrow\ z=-1\text{ and }X_x\ge Y_y.
\]
There is one fixed permutation of the input cube, depending only on $k$,
with the following properties for either activation.
\begin{enumerate}
\item The target has a Boolean circuit of size $\Theta(k)$ and a
width-one teacher for the activation $\phi$, with every weight at most
one in absolute value and signed output at least one everywhere.
\item After at most $4k+3$ calibration examples, there are
$T_k=2^k-k-1\ge2^{k-1}$ designated rounds. At every fixed width
$N\ge1$, the posterior wrong-label probability at each of these rounds
tends to one as $\sigma\downarrow0$.
\item For every fixed $\gamma>1$ and $\tau>0$ in
\eqref{eq:scaling}, the same wrong-label probabilities tend to one as
$N\to\infty$.
\item For every $\eta\in(0,1)$ there is $\tau_0(k,\eta)>0$ such that,
for every $0<\tau\le\tau_0$, every $\gamma\ge1$, and every $N\ge1$,
posterior sampling makes all $T_k$ designated mistakes with probability
at least $1-\eta$. The bound also makes each designated posterior
majority wrong, and violates the $2/3$ correctness criterion, when
$\eta$ is chosen sufficiently small.
\item For the same prior and target, the realizable IID sample complexity
is $O((d^2+\log(1/\beta))/\varepsilon)$ for correct-label posterior
probability greater than $2/3$ outside a set of input mass at most
$\varepsilon$, with sample confidence at least $1-\beta$.
This holds at each fixed width for sufficiently small $\sigma$, and
under the direct width limits specified above.
\end{enumerate}
No choice among tied predictors is used. Every network that would make
the correct next prediction is included in the energy comparison.
\end{theorem}

The discontinuity distinguishes this example from ordinary tanh or ReLU.
The theorem keeps the Gaussian prior and its usual layer normalizations
unchanged, but changes the activation. The small-prefactor assertion is
not a claim about the unit-prefactor case $\gamma=1,\tau=1$.

\section{An exact minimum over every hidden width}

We first prove the simpler statement for $\phi=\phi_{\rm lin}$.
Set $b_j=Nv_j$. The Gaussian rate energy and output become
\begin{equation}\label{eq:energy}
 E_N=\frac dN\sum_j\|u_j\|^2+\frac1N\sum_jb_j^2,
 \qquad F(x)=\frac1N\sum_jb_j\phi(u_j\cdot\bar x).
\end{equation}
Consider any finite linearly separable labeled set, and put
$z_i=y_i\bar x_i$. Let $a_*$ minimize $\|a\|^2$ under
$z_i\cdot a\ge1$, and write $Q=\|a_*\|^2>0$.
The minimizer is unique. The normal-cone condition for this Euclidean
projection gives active constraints and nonnegative multipliers with
\[
 a_*=\sum_i\mu_i z_i,\qquad
 \mu_i>0\Longrightarrow z_i\cdot a_*=1,\qquad
 \sum_i\mu_i=Q.
\]
The last identity follows by taking the inner product with $a_*$.
Put $\rho_i=\mu_i/Q$.

\begin{theorem}[Exact all-width affine reduction]\label{thm:energy}
At every $N\ge1$, the minimum of \eqref{eq:energy} over all networks
satisfying $y_iF(x_i)\ge1/2$ is
\begin{equation}\label{eq:profile}
 \Psi_d(Q)=\min_{r\ge0}\left\{dQr^2+\frac1{4(1+r)^2}\right\}.
\end{equation}
Its unique minimizing radius is strictly positive and obeys
\[
 4dQr(1+r)^3=1.
\]
The minimum is attained by cloning
$u_j=r a_*$ and $b_j=1/[2(1+r)]$.
The function $\Psi_d(Q)$ is strictly increasing in $Q>0$ and is always
strictly less than $1/4$.
\end{theorem}
\begin{proof}
Define the active-data average
\[
 S(u)=\sum_i\rho_i\phi(z_i\cdot u)
     =\sum_i\rho_i\sgn(z_i\cdot u)+a_*\cdot u/Q.
\]
Since the sign average has absolute value at most one,
\[
 |S(u)|\le1+\|u\|/\sqrt Q.
\]
Let $\E_j$ mean uniform averaging over the $N$ neurons and set
$r^2=\E_j\|u_j\|^2/Q$. Oddness of $\phi$ and the training margins give
\begin{align*}
 \frac12
 &\le\E_j[b_jS(u_j)]\\
 &\le\sqrt{\E_jb_j^2}\,
       \sqrt{\E_j(1+\|u_j\|/\sqrt Q)^2}
 \le\sqrt{\E_jb_j^2}\,(1+r).
\end{align*}
The last step uses $\E_j\|u_j\|\le\sqrt{\E_j\|u_j\|^2}$.
Therefore $E_N\ge dQr^2+1/[4(1+r)^2]$.
The displayed clones attain this bound: every active signed
preactivation is $r>0$, and every other one is at least $r$.

The scalar objective has derivative $2dQr-1/[2(1+r)^3]$ and strictly
positive second derivative. Its derivative at zero is negative, so its
unique minimum is at $r>0$ and below its value $1/4$ at zero.
For $Q_1>Q_0$, evaluating the profile for $Q_0$ at the positive minimizer
$r_1$ for $Q_1$ gives
\[
 \Psi_d(Q_1)-\Psi_d(Q_0)\ge d(Q_1-Q_0)r_1^2>0.
\]
\end{proof}

The theorem is an inequality over all neurons and all nonlinear output
functions. It does not restrict competitors to cloned networks.
For a threshold truth table $g$, apply it to the complete labeled cube:
if $Q(g)$ is the minimum affine unit-margin squared norm, then the
strict-margin cell has energy infimum $\Psi_d(Q(g))$.
A small common readout amplification makes the cloned margins strict.

\begin{lemma}[The nonthreshold gap for the linear residual]
\label{lem:nonthreshold}
If a Boolean function $g$ is not an affine threshold function, then every
network satisfying $g(x)F(x)\ge1/2$ on the full cube has $E_N\ge1/4$.
\end{lemma}
\begin{proof}
Put $z_x=g(x)\bar x$. Since these finitely many vectors do not admit a
strict separating linear functional, the origin belongs to their convex
hull. Indeed otherwise strict separation of the origin from that compact
convex hull would give an affine threshold representation of $g$.
Choose probabilities $\rho_x$ with $\sum_x\rho_x z_x=0$. Then
\[
 S(u):=\sum_x\rho_x\phi(z_x\cdot u)
      =\sum_x\rho_x\sgn(z_x\cdot u),\qquad |S(u)|\le1.
\]
Averaging the margins gives
$1/2\le\E_j[b_jS(u_j)]\le\E_j|b_j|\le\sqrt{\E_jb_j^2}$.
The readout part of the energy alone is therefore at least $1/4$.
\end{proof}

\section{The same exact reduction with a bounded residual}

The bounded activation gives the same online separation without an
unbounded output from a hidden neuron. We prove a slightly more general
statement for
\[
 \phi_c(t)=\sgn(t)+c\tanh(t),\qquad 0<c\le1.
\]
The coefficient $c$ is fixed across input dimensions and learner widths.

\begin{theorem}[A bounded threshold residual]\label{thm:bounded}
For any nonempty linearly separable labeled subset of the Boolean cube,
let $a_*,Q$ be its affine unit-margin minimizer and squared norm as above.
At every width, the minimum training-margin energy for $\phi_c$ is
\begin{equation}\label{eq:bounded-profile}
 \Xi_{d,c}(Q)=\min_{r\ge0}
       \left\{dQr^2+\frac1{4(1+c\tanh r)^2}\right\}.
\end{equation}
The unique minimizing radius is positive and satisfies
\begin{equation}\label{eq:bounded-stationary}
 4dQr(1+c\tanh r)^3=c\operatorname{sech}^2r.
\end{equation}
Cloning $u_j=r a_*$ and $b_j=1/[2(1+c\tanh r)]$ attains the minimum.
The profile is strictly increasing in $Q$ and is strictly below $1/4$.
Every globally nonthreshold Boolean output still has energy at least
$1/4$.
\end{theorem}
\begin{proof}
Every signed augmented input has norm $\sqrt{d+1}$. The unit-margin
constraint at any training point implies $Q\ge1/(d+1)$, so $dQ\ge1/2$.
The scalar objective in~\eqref{eq:bounded-profile} has negative derivative
at zero and strictly increasing derivative: its negative derivative
term is proportional to
$c\operatorname{sech}^2r/(1+c\tanh r)^3$, which decreases strictly.
Thus its minimizer is unique and positive and satisfies
\eqref{eq:bounded-stationary}. In particular
$r\le c/(4dQ)\le1/2$.

Write
\[
 h=1+c\tanh r,\quad
 \alpha=1+c(\tanh r-r\operatorname{sech}^2r),\quad
 \beta=c\operatorname{sech}^2r,\quad
 \lambda=h^{-2},\quad K=\lambda^2/4.
\]
Then $h=\alpha+\beta r$, $\alpha\ge1$, $0<\beta\le1$, and $K\le1/4$.
For any score $s\in[-\sqrt2,\sqrt2]$,
\begin{equation}\label{eq:bounded-tangent}
 -\alpha+\beta s\le\phi_c(s)\le\alpha+\beta s.
\end{equation}
For $s>0$, the upper inequality is the concave tangent bound for tanh
at $r$. For $s<0$, $\phi_c(s)\le-1$, whereas
$\alpha+\beta s\ge1-\sqrt2>-1$. At zero it is immediate. Oddness gives
the lower inequality. The fact $\alpha\ge1$ follows from
$\tanh r-r\operatorname{sech}^2r\ge0$, whose derivative is
$2r\operatorname{sech}^2r\tanh r\ge0$.

Use the active KKT probabilities $\rho_i$ and put
$S(u)=\sum_i\rho_i\phi_c(z_i\cdot u)$.
For any atom and any readout $b$, completing the square gives
\[
 d\|u\|^2+b^2-\lambda bS(u)
 \ge d\|u\|^2-KS(u)^2.
\]
If $\|u\|\le1/\sqrt d$, every score has magnitude at most $\sqrt2$.
Equation~\eqref{eq:bounded-tangent} and the KKT mean give
\[
 |S(u)|\le\alpha+\beta\|u\|/\sqrt Q.
\]
Writing $t=\|u\|/\sqrt Q$, the atomwise lower bound is therefore at least
\[
 \min_{t\ge0}\{dQt^2-K(\alpha+\beta t)^2\}
 =dQr^2-Kh^2=:\zeta.
\]
Indeed $dQ-K\beta^2\ge1/4>0$, and stationarity says
$dQr=K\beta h$, making $t=r$ its unique minimum. It also gives
$\zeta=-\alpha/(4h^3)<0$.
If $\|u\|\ge1/\sqrt d$, boundedness gives instead
\[
 d\|u\|^2-KS(u)^2\ge1-4K\ge0>\zeta.
\]
Thus the same floor $\zeta$ holds for every atom. Averaging and using
the training margins yields
\[
 E_N\ge\lambda/2+\zeta
       =dQr^2+\frac1{4h^2}.
\]
The proposed clones attain equality on the active constraints and satisfy
all others by monotonicity of $\phi_c$ on positive scores. Strict
monotonicity in $Q$ follows by evaluating the profile at the positive
minimizer for the larger $Q$, as in Theorem~\ref{thm:energy}. Its value is
below $1/4$ because its derivative at zero is negative.

For the final assertion, choose the zero-mean convex-hull probabilities
for a nonthreshold function as in Lemma~\ref{lem:nonthreshold}.
The tangent estimate~\eqref{eq:bounded-tangent} at $r=0$ has
$\alpha=1,\beta=c$. Thus $|S(u)|\le1$ on $\|u\|\le1/\sqrt d$.
On the remaining atoms $|S(u)|\le1+c\le2$. With square-completion
coefficient one,
\[
 d\|u\|^2+b^2-bS(u)\ge d\|u\|^2-S(u)^2/4\ge-1/4
\]
in both regions; in the large-atom region the right side is actually
nonnegative. Averaging and using $\E_j b_jS(u_j)\ge1/2$ proves
$E_N\ge1/4$.
\end{proof}

From now on write $J_d(Q)$ for $\Psi_d(Q)$ with the linear residual,
or $\Xi_{d,1}(Q)$ with the bounded residual. Both profiles are strictly
increasing and below $1/4$, and both have explicit cloned minimizers.
The same arguments also cover every fixed $0<c\le1$ in the bounded
case, with $J_d=\Xi_{d,c}$.

\section{A fixed input order with exponentially many affine mistakes}

We give the complete affine construction used by the network theorem.
Keep the target $f_k$ from Theorem~\ref{thm:main}, and abbreviate the two
encoded integers by $X$ and $Y$. The calibration examples all have
selector $z=-1$.
\begin{enumerate}
\item Take diagonal inputs $x=y$ equal to the all-positive vector,
each of its $k$ one-coordinate flips, and the negatives of these vectors.
Every label is positive.
\item For each $i\in\{0,\ldots,k-1\}$, take two negative examples with
$x_i=-1,y_i=+1$ and $x_j=+1,y_j=-1$ for $j<i$. For $j>i$, set the equal
bits $x_j=y_j$ all positive in one example and all negative in the other.
Their highest differing bit is $i$, so $X<Y$.
\end{enumerate}
Remove duplicate examples. The resulting calibration has at most $4k+2$
inputs. Let $q_t$ denote the input $z=+1$, $Y=0$, $X=t$; its target label
is negative for every $0\le t\le2^k-1$.

\begin{lemma}\label{lem:selector}
For $k\le t\le2^k-1$, the unique minimum-norm affine unit-margin
coefficient vector fitting the calibration and the negative example
$q_t$ is the vector of the affine score
\begin{equation}\label{eq:selector-score}
 a_t\cdot\bar x=-t-(t+1)z+\sum_{i=0}^{k-1}2^i(x_i-y_i).
\end{equation}
It fits all earlier $q_j$ with $j\le t$, and its absolute score is at
least one on the entire cube. Its squared norm is
\begin{equation}\label{eq:Qt}
 Q_t=t^2+(t+1)^2+2\sum_{i=0}^{k-1}4^i.
\end{equation}
For $t<2^k-1$, it predicts $q_{t+1}$ positively, and adding the correct
negative label raises the convex affine minimum by
\begin{equation}\label{eq:Qgap}
 Q_{t+1}-Q_t=4(t+1)>0.
\end{equation}
\end{lemma}
\begin{proof}
Write an arbitrary affine score as
\[
 b+cz+\sum_i a_i x_i+\sum_i h_i y_i,
\]
and set
\[
 u=b-c,\quad v=b+c,\quad
 w_i=(a_i-h_i)/2,\quad s_i=a_i+h_i,\quad U=u-1.
\]
The squared coefficient norm is
\begin{equation}\label{eq:affine-norm}
 \frac{u^2+v^2}{2}+2\sum_iw_i^2+\frac12\sum_i s_i^2.
\end{equation}
The paired positive diagonal constraints give
\[
 U\ge0,\qquad \left|\sum_j s_j\right|\le U,\qquad
 \left|\sum_j s_j-2s_i\right|\le U.
\]
Subtracting the expressions inside the last two absolute values shows
$|s_i|\le U$, hence $\sum_i|s_i|\le kU$.
Average the two negative carry constraints at coordinate $i$. The upper
equal-bit terms cancel and leave
\[
 u-2w_i+2\sum_{j<i}w_j\le-1.
\]
Induction yields
\begin{equation}\label{eq:binary-lower}
 w_i\ge\sum_{j<i}w_j+\frac{u+1}{2}
       \ge2^i(1+U/2).
\end{equation}
At $q_t$, the non-selector part of the input contribution is
$2\sum_{\operatorname{bit}_i(t)=1}w_i-
\sum_{\operatorname{bit}_i(t)=0}s_i$. Its negative unit-margin constraint
therefore implies
\[
 v\le-1-2\sum_{\operatorname{bit}_i(t)=1}w_i
            +\sum_{\operatorname{bit}_i(t)=0}s_i
   \le-(2t+1)+(k-t)U\le-(2t+1),
\]
since $t\ge k$. Combining this with $u\ge1$ and
\eqref{eq:binary-lower} in \eqref{eq:affine-norm} gives
\[
 \|a\|^2\ge\frac{1+(2t+1)^2}{2}+2\sum_i4^i=Q_t.
\]
Equality requires $u=1$, $s_i=0$, $w_i=2^i$ and $v=-(2t+1)$,
which is exactly \eqref{eq:selector-score}.

Conversely the displayed score equals $1+2(X-Y)$ on $z=-1$, and
$-(2t+1)+2(X-Y)$ on $z=+1$. It fits every calibration label, and its score
at $q_j$ is $2j-2t-1\le-1$ for $j\le t$. Both branch scores are odd
integers on the full cube, so the absolute margin is at least one.
At $q_{t+1}$ the score is $+1$. Applying the same minimum calculation
to the augmented history gives $Q_{t+1}$, and subtraction gives
\eqref{eq:Qgap}.
\end{proof}

The input permutation starts with the calibration, then $q_k$, then
$q_{k+1},\ldots,q_{2^k-1}$, and finally every remaining input in
lexicographic order. These blocks are disjoint. Its calibration prefix
has length at most $4k+3$, and there are
$T_k=2^k-k-1\ge2^{k-1}$ designated queries. Just before each one, the
wrong affine witness has the exact convex margin minimum for all data
seen so far.

The target has a linear-size Boolean circuit. With $0/1$ bit notation,
comparison of the lowest $i+1$ bits obeys
\[
 G_{-1}=1,\qquad
 G_i=(X_i\wedge\neg Y_i)\ \vee\
       ((X_i\leftrightarrow Y_i)\wedge G_{i-1}).
\]
Each step uses a constant number of bounded-fan-in gates; conjoin the
last result with $z=-1$. Every input bit is essential: for $x_i$ use
$Y=2^i$ and all other $X$ bits zero; for $y_i$ use $X=0$ and all other
$Y$ bits zero; for $z$ take $X=Y$. A fan-in-two circuit with $d$ essential
inputs has at least $d-1$ gates, by counting edges in the connected
ancestry of its output. This proves the matching $\Theta(k)$ bound.


\section{The posterior comparison, including growing width}

The following estimate handles the discontinuity by using a neighborhood
whose preactivations stay away from zero.

\begin{lemma}[A product-box Gaussian comparison]\label{lem:box}
Suppose a fixed finite history has a cloned single-neuron witness of
energy $e_0$ with a full-cube nonzero affine preactivation at every input,
and whose output has the desired global weak margin and history labels.
Suppose every network in a competing event has energy at least $e_1>e_0$,
at every width. Then there are constants $\delta>0$ and $C\ge0$,
independent of $N$ and $\sigma$, such that the competing posterior mass
is at most
\begin{equation}\label{eq:posterior-bound}
 \exp\left[-\frac{\delta}{2\sigma^2}+CN\right]
\end{equation}
whenever $(d+2)N\sigma^2<e_1$.
The posterior is conditioned on the history and global margin.
\end{lemma}
\begin{proof}
In physical neuron coordinates
\[
 \theta_j=(\sqrt d\,u_j,b_j)\in\R^q,\qquad q=d+2,
\]
the rows are independent centered Gaussians with covariance
$s^2I_q$, where $s^2=N\sigma^2$, and
$E_N=N^{-1}\sum_j\|\theta_j\|^2$.
Amplify the witness's readout slightly, so every margin is strict and its
energy remains below $e_1$. Its preactivations are nonzero on the finite
cube. Consequently there is a fixed open coordinate box
$B\subset\R^q$ of volume $V>0$ about this amplified neuron such that
every neuron in $B$ has the same strict full-cube sign function and
individual signed output greater than $1/2$, and
\[
 R_B:=\sup_{\theta\in B}\|\theta\|^2<e_1.
\]
This uses continuity only within a region of fixed nonzero preactivation
signs. If every neuron lies in $B$, their average satisfies the same
strict margins and history. Thus the posterior denominator is at least
\[
 \left[V(2\pi s^2)^{-q/2}e^{-R_B/(2s^2)}\right]^N.
\]
For the numerator, $NE_N/s^2$ is chi-square with $Nq$ degrees of freedom.
The exponential-moment bound, valid for $e_1>qs^2$, gives
\[
 \log\Pp(E_N\ge e_1)
 \le-\frac{Ne_1}{2s^2}
       +\frac{Nq}{2}\left[1+\log\frac{e_1}{qs^2}\right].
\]
Dividing by the denominator cancels all logarithmic $s^2$ terms. With
$\delta=e_1-R_B>0$, the logarithm of the ratio is at most
\[
 -\frac{\delta}{2\sigma^2}
 +N\left\{\frac q2\left[1+\log\frac{2\pi e_1}{q}\right]-\log V\right\}.
\]
Replace the expression in braces by its maximum with zero to obtain
\eqref{eq:posterior-bound}.
\end{proof}

Before the designated query $q_{t+1}$ in Lemma~\ref{lem:selector}, the
cloned witness has energy $e_0=J_d(Q_t)$, has no zero affine score on
the cube, and predicts that query positively. Every network making the
correct negative prediction there belongs to a global-margin cell and
therefore satisfies the enlarged training set with margins at least
$1/2$. Theorems~\ref{thm:energy} and~\ref{thm:bounded} bound its energy by
$e_1=J_d(Q_{t+1})>e_0$. Lemma~\ref{lem:box} applies directly.
There is no appeal to continuity of the activation at zero or to a
compactness argument across its jumps.

For each fixed $N$, \eqref{eq:posterior-bound} tends to zero as
$\sigma\downarrow0$. Under \eqref{eq:scaling}, it is
\[
 \exp[-\delta N^\gamma/(2\tau^2)+CN].
\]
For every fixed $\gamma>1$ and $\tau>0$ this tends to zero; the tail
condition in Lemma~\ref{lem:box} also holds eventually.
For $\gamma\ge1$ and sufficiently small $\tau$, both the tail condition
and any desired bound on the competing probability hold at \emph{every}
$N\ge1$: require
\[
 \tau^2<\frac{e_1}{d+2},\qquad
 \frac{\delta}{2\tau^2}-C\ge\log(T_k/\eta).
\]
Taking common constants over the finite list of designated histories
and using the union bound proves all online assertions of
Theorem~\ref{thm:main}. Independence between posterior samples at different
rounds is unnecessary.

\section{A width-one teacher and polynomial IID learning}

The final affine witness $a_{2^k-1}$ represents the target on the entire
cube: on $z=-1$ its score is $1+2(X-Y)$, and on $z=+1$ it is always
negative. Multiply all its affine coefficients by a sufficiently small
positive number so that every one has magnitude at most one, and use
readout weight one. All preactivations remain nonzero and retain their
signs. Both activations have $|\phi(t)|>1$ for $t\ne0$, so the resulting
width-one teacher has signed output greater than one. The linear Boolean circuit
and its matching lower bound were given with the selector construction.

\begin{theorem}[IID guarantee for the same prior]\label{thm:iid}
Fix any affine threshold target $f$ on $\{-1,1\}^d$, any input
distribution $\mathcal D$, and $\varepsilon,\beta\in(0,1)$. If
\[
 m\ge\frac4\varepsilon
       \bigl[(d+1)^2\log2+\log(1/\beta)\bigr],
\]
then with probability at least $1-\beta$ over $m$ IID examples labeled
by $f$, the posterior correct-label probability exceeds $2/3$ outside a
set of $\mathcal D$-mass at most $\varepsilon$ in each of these regimes:
\begin{enumerate}
\item at every fixed width, for all sufficiently small $\sigma>0$;
\item under \eqref{eq:scaling}, at every sufficiently large width for
any fixed $\gamma>1$ and $\tau>0$;
\item at every width and every $\gamma\ge1$ in \eqref{eq:scaling},
when $\tau>0$ is sufficiently small.
\end{enumerate}
The variance and width thresholds can be chosen independently of the
training sample and $\mathcal D$.
\end{theorem}
\begin{proof}
The full target has a cloned witness of energy
$e_f=J_d(Q(f))<1/4$, with nonzero full-cube affine scores.
Lemma~\ref{lem:nonthreshold} and Theorem~\ref{thm:bounded} put every
nonthreshold output cell in the
event $E_N\ge1/4$. The full-target witness is consistent with every
sample labeled by $f$. Lemma~\ref{lem:box}, with $e_1=1/4$, therefore
bounds the total posterior mass of all nonthreshold functions uniformly
over every such history. This bound tends to zero in the first two
regimes and can be made at most $\varepsilon/12$ uniformly in the third.
No count of nonthreshold functions is needed.

Let $\mathcal H_d$ be the threshold truth tables on the cube. The
$2^d$ hyperplanes in the $(d+1)$-dimensional coefficient space give
\[
 |\mathcal H_d|\le\sum_{j=0}^{d+1}\binom{2^d}{j}
 \le(2^d+1)^{d+1}\le2^{(d+1)^2}.
\]
A consistency union bound shows that, with probability at least $1-\beta$,
every sample-consistent threshold function has population error at most
$\varepsilon/4$. On this sample event, posterior mean population error
is at most $\varepsilon/4+\varepsilon/12=\varepsilon/3$ once the
nonthreshold mass is small enough. An input with correct-label posterior
probability at most $2/3$ has wrong-label probability at least $1/3$.
Markov's inequality bounds the mass of those inputs by $\varepsilon$.
\end{proof}

\section{What the activation contributes}

The threshold component supplies a signed output of magnitude one even
when its incoming weights are arbitrarily small. Either residual
then makes a small positive incoming scale strictly preferable to zero:
its energy optimum is below $1/4$ and depends strictly on the affine
margin norm. This creates an energy order at every width while keeping
a width-one bounded teacher for every affine threshold target.

The elementary linear-residual calculation also works for every fixed
$\sgn(t)+ct$ with $c>0$, replacing $1+r$ by $1+cr$ in its profile.
For the bounded family, Theorem~\ref{thm:bounded} already permits every
fixed $0<c\le1$. Its uniform distance from pure sign is at most $c$.
The residual may therefore be arbitrarily small, but its coefficient is
fixed across widths and dimensions. At $c=0$ the strict energy order
disappears, so that limiting case needs a different argument.

This example uses an unchanged Gaussian prior and actual posterior
conditioning, not a favorable choice of a minimum-energy predictor.
Its activation is discontinuous. Replacing the jump by a fixed continuous
transition is a substantive mathematical change, and the present proof
does not justify that replacement. The question for ordinary tanh,
ReLU, or hardtanh under this same prior remains separate.

\end{document}
